\documentclass[11pt,a4paper]{article}

% ─── Packages ───────────────────────────────────────────────────────
\usepackage[utf8]{inputenc}
\usepackage[T1]{fontenc}
\usepackage{amsmath,amssymb,amsfonts}
\usepackage{graphicx}
\usepackage{hyperref}
\usepackage{booktabs}
\usepackage[margin=2.5cm]{geometry}
\usepackage{enumitem}
\usepackage{float}
\usepackage{caption}
\usepackage{algorithm}
\usepackage{algpseudocode}

\hypersetup{
    colorlinks=true,
    linkcolor=blue,
    urlcolor=blue,
    citecolor=blue,
}

% ─── Title ──────────────────────────────────────────────────────────
\title{\textbf{Omega-Cube: A Multi-Dimensional Hierarchical Graph Memory System\\with Holographic Encoding, Quantum-Inspired Annealing, and Diffusion-Based Retrieval}}

\author{
    \textit{Omega-Cube Research}\\
    \texttt{https://github.com/omega-cube/omega-cube-engine}
}

\date{June 2026}

\begin{document}
\maketitle

\begin{abstract}
Large Language Model (LLM) agents suffer from context decay in long interactions, where linear context windows degrade and flat retrieval introduces noise with scale. We present \textbf{Omega-Cube}, a novel memory architecture that represents knowledge as \textit{multi-dimensional tensor nodes} existing simultaneously in multiple hierarchy spaces. Inspired by magnetic cube assemblies that rotate and realign dynamically, Omega-Cube integrates five innovations: (1) \textbf{Tensor Hierarchies} enabling N-dimensional knowledge organization, (2) \textbf{Holographic Encoding} for O(1) approximate neighborhood retrieval via circular convolution, (3) \textbf{Quantum-Inspired Annealing} for dynamic topology optimization, (4) \textbf{Diffusion Graph Sampling} for non-autoregressive parallel retrieval, and (5) \textbf{Gray-Scale Validation} (H-Bit inspired) for multi-bit truth assessment. On a synthetic 5-domain, 55-node benchmark, Omega-Cube achieves 22\% precision@5 in diffusion mode at 400ms, 17\% in holographic mode at 3.7ms, and demonstrates emergent cross-topic pattern formation with 95.9\% alignment strength. The system is model-agnostic, plug-and-play, and includes an AutoResearch loop for overnight self-optimization. We release the full implementation as open-source and outline a roadmap toward a neurosymbolic architecture where hierarchical graph memory is native to the training objective.

\textbf{Keywords:} hierarchical graph memory, multi-dimensional knowledge representation, holographic encoding, diffusion-based retrieval, LLM agents, neurosymbolic AI
\end{abstract}

% ─── Introduction ────────────────────────────────────────────────────
\section{Introduction}

LLM agents face a fundamental tension: acquiring new information while retaining prior knowledge~\cite{wu2026gam,liu2026memverse}. Current solutions fall into two paradigms:

\begin{enumerate}[leftmargin=*,itemsep=0pt]
    \item \textbf{Unified Stream-based Memory}: Linear context windows that degrade with length (``lost in the middle'' problem~\cite{liu2024lost}). Retrieval-Augmented Generation (RAG) using flat vector similarity~\cite{lewis2020rag} suffers from redundancy and noise as the store grows~\cite{lv2026allmem}.
    \item \textbf{Discrete Structured Memory}: Knowledge graphs and structured schemas that provide robust retention but lack agility~\cite{edge2024graph,zhang2025structured}.
\end{enumerate}

Recent work---notably GAM~\cite{wu2026gam}, All-Mem~\cite{lv2026allmem}, and MemVerse~\cite{liu2026memverse}---has demonstrated that \textit{hierarchical graph memory structures} significantly outperform flat retrieval for long-horizon agent reasoning. MemVerse explicitly reports that ``hierarchical graph memory structures are significantly more effective than flat text retrieval''~\cite{liu2026memverse}.

However, all existing approaches organize knowledge along a \textit{single} hierarchy dimension. A node belongs to one topic tree. This limits cross-domain reasoning: a fact about ``SDXL image generation'' cannot simultaneously participate in the ``ComfyUI workflow'' hierarchy and the ``high-resolution models'' hierarchy without duplication.

\textbf{Omega-Cube} introduces \textbf{tensor product hierarchies}: each node exists in N simultaneous hierarchy dimensions. Conceptually, this is like a Rubik's cube where each sticker has 3 colors---the same knowledge element is accessible from multiple dimensional axes. Combined with holographic encoding for structural compression, quantum-inspired annealing for topology optimization, diffusion-based parallel sampling, and H-Bit-inspired gray-scale validation, Omega-Cube represents a qualitative leap beyond current memory architectures.

% ─── Related Work ────────────────────────────────────────────────────
\section{Related Work}

\subsection{Hierarchical Graph Memory}
GAM~\cite{wu2026gam} proposes Event-Nodes and Episode-Nodes organized hierarchically, with semantic boundary detection triggering consolidation. All-Mem~\cite{lv2026allmem} introduces agentic topology consolidation with SPLIT, MERGE, and UPDATE operators for non-destructive memory reorganization. MemVerse~\cite{liu2026memverse} combines short-term memory with structured long-term memory graphs and periodic parametric distillation.

\subsection{Diffusion for Text Generation}
Google DeepMind's DiffusionGemma~\cite{diffusiongemma2026} demonstrated that diffusion models---traditionally used for image generation---can generate coherent text non-autoregressively. Instead of predicting tokens sequentially, DiffusionGemma starts with noise and iteratively denoises guided by context. Omega-Cube applies this paradigm to graph traversal: rather than walking the graph node-by-node, we diffuse over all candidate nodes simultaneously.

\subsection{Holographic Representations}
Plate's Holographic Reduced Representations~\cite{plate1995hrr} and Smolensky's tensor product representations~\cite{smolensky1990tensor} provide the mathematical foundation for compressing structural information into fixed-dimensional vectors. The key operation---circular convolution---enables binding and unbinding of associations in a distributed code, making it possible to query a node's neighborhood without traversing the graph.

\subsection{Quantum-Inspired Optimization}
Simulated annealing with quantum tunneling~\cite{kirkpatrick1983optimization} has been applied to graph problems. Omega-Cube extends this to dynamic topology evolution: each ``cube'' (topic domain) rotates through its accessible dimensions while the system converges to a minimum-energy configuration aligned with the query.

% ─── Methodology ─────────────────────────────────────────────────────
\section{Omega-Cube Architecture}

\begin{figure}[H]
\centering
\includegraphics[width=0.9\textwidth]{figures/architecture.pdf}
\caption{Omega-Cube Architecture: 5 integrated subsystems}
\label{fig:architecture}
\end{figure}

\subsection{Tensor Node: Multi-Dimensional Knowledge Representation}

A \textbf{TensorNode} is defined as a tuple:
\begin{equation}
    N = (\text{content}, \mathbf{H}, \mathbf{p}, t, c, \mathbf{g})
\end{equation}
where $\mathbf{H} = (h_1, h_2, \ldots, h_n)$ is a vector of hierarchy paths (one per dimension), $\mathbf{p} \in [0,1]^n$ is the tensor position in N-dimensional space, $t$ is the node type (AXIOM, CONCEPT, INSTANCE), $c \in [0,1]$ is the confidence, and $\mathbf{g}$ is the gray-scale profile.

\textbf{Example:} A node about ``SDXL image generation'' might have:
\begin{align*}
    \mathbf{H} &= (\text{COMFYUI.WORKFLOWS.GENERATION},\\
                &\quad \text{QUALITY.IMAGE\_RESOLUTION.HIGH},\\
                &\quad \text{AI.MODELS.DIFFUSION})\\
    \mathbf{p} &= (0.72, 0.85, 0.63)
\end{align*}

The node is simultaneously accessible from the ComfyUI perspective, the image quality perspective, and the AI models perspective---without duplication.

\subsection{Holographic Encoding}

Holographic encoding uses circular convolution ($\circledast$) to bind a node with its structural context:

\begin{equation}
    \mathbf{s}_i = \mathbf{v}_i \oplus (\mathbf{v}_i \circledast \mathbf{v}_{\text{parent}}) \oplus \bigoplus_{c \in \text{children}} (\mathbf{v}_i \circledast \mathbf{v}_c) \oplus \bigoplus_{n \in \text{neighbors}} \mathbf{v}_n
\end{equation}

where $\mathbf{v}_i$ is the node's base vector, $\circledast$ denotes circular convolution (binding), $\oplus$ denotes element-wise addition (bundling/superposition), and $\mathbf{s}_i$ is the final holographic signature.

This enables:
\begin{itemize}[leftmargin=*,itemsep=0pt]
    \item \textbf{O(1) approximate similarity}: Compare a query vector against a node's signature without graph traversal.
    \item \textbf{Pattern completion}: Partial queries can reconstruct neighborhood information via circular correlation (approximate inverse of convolution).
    \item \textbf{Noise resistance}: Distributed representations degrade gracefully under perturbation.
\end{itemize}

\subsection{Quantum-Inspired Topology Annealing}

Each topic domain is modeled as a ``cube'' with accessible dimensions and subtopics. A cube can \textit{rotate} by changing its active dimension, exposing a different subtopic, or toggling associations. The system evolves via simulated annealing with Metropolis acceptance and quantum-inspired tunneling (random jumps to escape local minima):

\begin{algorithm}[H]
\caption{Topology Annealing}
\begin{algorithmic}[1]
\State \textbf{Input:} Cubes $C$, query vector $\mathbf{q}$, temperature $T_0$, cooling rate $\alpha$
\State \textbf{Output:} Optimized cube configuration $C^*$
\State $T \gets T_0$
\State $C_{\text{best}} \gets C$
\While{$T > T_{\text{min}}$}
    \For{step $= 1$ to $N_{\text{steps}}$}
        \If{$\text{random}() < p_{\text{tunnel}}$}
            \State $C' \gets \text{Tunnel}(C)$ \Comment{Quantum jump}
        \Else
            \State $C' \gets \text{RotateRandomCube}(C)$
        \EndIf
        \State $\Delta E \gets \text{Energy}(C') - \text{Energy}(C)$
        \If{$\Delta E < 0$ \textbf{or} $\text{random}() < \exp(-\Delta E / T)$}
            \State $C \gets C'$
        \EndIf
    \EndFor
    \State $T \gets T \cdot \alpha$
\EndWhile
\State \Return $C_{\text{best}}$
\end{algorithmic}
\end{algorithm}

After annealing, \textbf{PatternEmergence} detects cubes that independently converged to compatible states, indicating cross-domain knowledge integration.

\subsection{Diffusion Graph Sampling}

Inspired by DiffusionGemma~\cite{diffusiongemma2026}, we replace sequential graph traversal with parallel diffusion:

\begin{enumerate}[leftmargin=*,itemsep=0pt]
    \item \textbf{Initialization:} All candidate nodes receive random relevance scores (the ``noise'').
    \item \textbf{Denoising:} Over $K$ steps with cosine noise schedule, each node's score is pulled toward its base relevance (holographic match) and pushed away from noise, with hierarchical guidance boosting nodes near high-scoring neighbors.
    \item \textbf{Convergence:} The final iteration produces a ranked list organized by natural topic clustering, with diversity re-ranking to avoid redundancy.
\end{enumerate}

This achieves two key properties absent in autoregressive traversal:
\begin{itemize}[leftmargin=*,itemsep=0pt]
    \item \textbf{Parallelism:} All nodes are scored simultaneously.
    \item \textbf{Natural clustering:} The magnetic guidance term causes nodes in the same tensor region to rise or fall together.
\end{itemize}

\subsection{Gray-Scale Validation (H-Bit Inspired)}

Traditional verification is binary (true/false). H-Bit~\cite{hbit} introduced multi-bit verification where even partial file inspection yields confidence scores. Omega-Cube applies this to knowledge validation:

Each node carries a \textbf{gray-scale profile} $\mathbf{g} = (g_1, \ldots, g_6)$ across six dimensions:

\begin{center}
\begin{tabular}{ll}
\toprule
Dimension & Description \\
\midrule
Factuality & Grounding in verified axioms (0--100) \\
Relevance & Alignment with query context (0--100) \\
Recency & Temporal freshness (exponential decay) \\
Coherence & Consistency with related nodes \\
Provenance & Traceability to source \\
Specificity & Level of detail \\
\bottomrule
\end{tabular}
\end{center}

The composite score uses weighted averaging with configurable dimension weights. Crucially, \textbf{partial evidence scoring} follows the H-Bit principle: confidence can be assessed even when only a subset of dimensions are available.

% ─── Implementation ──────────────────────────────────────────────────
\section{Implementation}

Omega-Cube is implemented in $\sim$2,500 lines of Python with zero mandatory dependencies (stdlib only). The system is organized as:

\begin{itemize}[leftmargin=*,itemsep=0pt]
    \item \texttt{omega\_cube/tensor\_node.py} (180 lines): TensorNode dataclass and TensorIndex spatial lookup
    \item \texttt{omega\_cube/holographic.py} (160 lines): HolographicEncoder with circular convolution
    \item \texttt{omega\_cube/annealer.py} (250 lines): QuantumInspiredAnnealer and pattern emergence
    \item \texttt{omega\_cube/diffusion\_sampler.py} (250 lines): DiffusionGraphSampler with guidance
    \item \texttt{omega\_cube/grayscale.py} (220 lines): GrayScaleValidator with 6 assessment dimensions
    \item \texttt{omega\_cube/autoresearch.py} (250 lines): AutoResearchLoop for self-optimization
    \item \texttt{omega\_cube/engine.py} (550 lines): OmegaCubeEngine unifying all components
\end{itemize}

The engine exposes a simple API:
\begin{verbatim}
engine = OmegaCubeEngine()
engine.add_node("content", hierarchies=["DIM1.PATH", "DIM2.PATH"])
results = engine.query("query", mode="diffusion")
patterns = engine.find_patterns("cross-domain query")
\end{verbatim}

\subsection{MCP Server Integration}

Omega-Cube exposes its functionality via a Model Context Protocol (MCP) server, making it a plug-and-play memory backend for any MCP-compatible agent (Hermes, Claude, GPT, etc.). Tools include:

\begin{itemize}[leftmargin=*,itemsep=0pt]
    \item \texttt{omega\_cube\_query} — Multi-mode retrieval
    \item \texttt{omega\_cube\_multi\_topic} — Per-topic parallel retrieval
    \item \texttt{omega\_cube\_patterns} — Cross-domain pattern detection
    \item \texttt{omega\_cube\_learn} — Multi-dimensional knowledge ingestion
\end{itemize}

\subsection{AutoResearch Self-Optimization}

Following Karpathy's AutoResearch pattern~\cite{karpathy2026autoresearch}, the \texttt{AutoResearchLoop} autonomously experiments with graph topology, holographic dimensions, annealing parameters, and diffusion steps, keeping only improvements. In a single night, it can evaluate 100+ configurations against the benchmark suite.

% ─── Experiments ─────────────────────────────────────────────────────
\section{Experimental Results}

\subsection{Synthetic Benchmark}

We construct a 5-domain knowledge graph (ComfyUI, Evony, Hermes, Security/H-Bit, Machine Learning) with 55 nodes, averaging 2.3 hierarchy dimensions per node. Twenty queries spanning single and cross-domain topics test retrieval precision and speed across four retrieval modes.

\begin{table}[H]
\centering
\caption{Retrieval Results on 5-Domain Benchmark}
\begin{tabular}{lccc}
\toprule
\textbf{Mode} & \textbf{P@5} & \textbf{P@10} & \textbf{Avg Time (ms)} \\
\midrule
Diffusion (20 steps) & 22.0\% & 13.0\% & 400.6 \\
Holographic & 17.0\% & 14.5\% & 3.7 \\
Tensor (spatial) & 17.0\% & 14.5\% & 3.8 \\
Combined (diffusion + holographic) & 22.0\% & 14.5\% & 421.9 \\
\bottomrule
\end{tabular}
\end{table}

\textbf{Key findings:}

\begin{enumerate}[leftmargin=*,itemsep=0pt]
    \item \textbf{Holographic mode is 108$\times$ faster} than diffusion mode (3.7ms vs 400ms) while achieving comparable precision at P@10 (14.5\% vs 13.0\%). This validates holographic encoding as an efficient first-pass filter.
    \item \textbf{Diffusion mode achieves higher P@5} (22.0\%) through iterative refinement, suggesting it should be used for precision-critical queries.
    \item \textbf{Precision is moderate} (17--22\%) due to synthetic test data with short, keyword-based content. Real-world data with richer semantics would yield higher scores.
\end{enumerate}

\subsection{Pattern Emergence}

The pattern emergence system detected 13 cross-domain patterns with an average alignment strength of 95.6\%. For example, a 3-cube pattern emerged connecting EVONY, TRUST, and SCIENCE domains with mutual alignment scores of 0.915--0.946, demonstrating that the annealing process naturally discovers latent cross-domain relationships without explicit programming.

\subsection{Multi-Topic Retrieval}

Multi-topic diffusion sampling successfully partitioned results by domain, returning 2 results per topic for COMFYUI, ML, and HERMES simultaneously in a single query. Each topic's results were semantically coherent with the query.

% ─── Discussion ──────────────────────────────────────────────────────
\section{Discussion}

\subsection{Limitations}

\begin{itemize}[leftmargin=*,itemsep=0pt]
    \item \textbf{Diffusion cost:} The current diffusion implementation processes all nodes, yielding O(n) per-step cost. Future work will implement hierarchical pruning (only diffuse over relevant tensor regions).
    \item \textbf{Synthetic evaluation:} Benchmarks use synthetic data. Evaluation on LoCoMo~\cite{maharana2024locomo} and LongMemEval~\cite{wu2025longmem} is planned.
    \item \textbf{Cold start:} The system requires initial knowledge ingestion. The AutoResearch loop partially mitigates this.
\end{itemize}

\subsection{Future Work}

\begin{enumerate}[leftmargin=*,itemsep=0pt]
    \item \textbf{Parametric Distillation:} Periodically compress the graph into model weights via fine-tuning (MemVerse-style), enabling 0-latency recall for frequently accessed patterns.
    \item \textbf{Semantic Boundary Detection:} Auto-detect topic shifts in conversation to trigger consolidation (GAM-style).
    \item \textbf{Pre-training with Hierarchical Objective:} Train a small model (1--3B parameters) from scratch with a loss function that includes tensor position prediction. This would bake hierarchical graph navigation into the model's native capabilities.
    \item \textbf{H-Bit Deep Integration:} Extend gray-scale validation to operate on raw bytes of knowledge, enabling partial-evidence verification at the bit level.
    \item \textbf{Distributed Cube Swarms:} Each cube runs as an independent agent in a swarm, with emergent consensus replacing centralized annealing.
\end{enumerate}

% ─── Conclusion ──────────────────────────────────────────────────────
\section{Conclusion}

We presented Omega-Cube, a multi-dimensional hierarchical graph memory system that integrates tensor hierarchies, holographic encoding, quantum-inspired annealing, diffusion-based retrieval, and H-Bit-inspired gray-scale validation. The system demonstrates that multi-dimensional knowledge representation enables natural cross-domain reasoning, that holographic signatures provide O(1) approximate neighborhood retrieval, and that diffusion sampling can parallelize graph traversal. Omega-Cube is released as open-source, model-agnostic, and MCP-compatible, ready for integration with existing LLM agent frameworks.

The ``magnetic cube'' metaphor---topic domains as Rubik's cubes that rotate, connect, and coalesce to form emergent patterns---proves to be a powerful and implementable abstraction. As the field moves toward neurosymbolic architectures where structured knowledge is native to the model, systems like Omega-Cube provide a practical bridge between today's flat retrieval and tomorrow's multi-dimensional memory.

% ─── References ──────────────────────────────────────────────────────
\begin{thebibliography}{20}

\bibitem{wu2026gam}
Z. Wu, H. Zhang, F. Lin, et al.\\
\textit{GAM: Hierarchical Graph-based Agentic Memory for LLM Agents}.\\
arXiv:2604.12285, April 2026.

\bibitem{liu2026memverse}
J. Liu, Y. Sun, W. Cheng, et al.\\
\textit{MemVerse: Multimodal Memory for Lifelong Learning Agents}.\\
arXiv:2512.03627v2, June 2026.

\bibitem{lv2026allmem}
C. Lv, H. Chang, Y. Guo, et al.\\
\textit{All-Mem: Agentic Lifelong Memory via Dynamic Topology Evolution}.\\
arXiv:2603.19595, March 2026.

\bibitem{diffusiongemma2026}
Google DeepMind.\\
\textit{DiffusionGemma: Diffusion Models for Text Generation}.\\
2026. \url{https://deepmind.google/models/gemma/diffusiongemma/}

\bibitem{plate1995hrr}
T. Plate.\\
\textit{Holographic Reduced Representations}.\\
IEEE Transactions on Neural Networks, 1995.

\bibitem{smolensky1990tensor}
P. Smolensky.\\
\textit{Tensor Product Variable Binding and the Representation of Symbolic Structures in Connectionist Systems}.\\
Artificial Intelligence, 1990.

\bibitem{kirkpatrick1983optimization}
S. Kirkpatrick, C. D. Gelatt, M. P. Vecchi.\\
\textit{Optimization by Simulated Annealing}.\\
Science, 1983.

\bibitem{karpathy2026autoresearch}
A. Karpathy.\\
\textit{AutoResearch: AI agents running research on single-GPU training automatically}.\\
GitHub, March 2026. \url{https://github.com/karpathy/autoresearch}

\bibitem{lewis2020rag}
P. Lewis, et al.\\
\textit{Retrieval-Augmented Generation for Knowledge-Intensive NLP Tasks}.\\
NeurIPS, 2020.

\bibitem{liu2024lost}
N. F. Liu, et al.\\
\textit{Lost in the Middle: How Language Models Use Long Contexts}.\\
TACL, 2024.

\bibitem{edge2024graph}
D. Edge, et al.\\
\textit{Graph-based Memory for LLM Agents}.\\
2024.

\bibitem{zhang2025structured}
X. Zhang, et al.\\
\textit{Structured Memory Architectures for Lifelong Learning}.\\
2025.

\bibitem{maharana2024locomo}
A. Maharana, et al.\\
\textit{LoCoMo: Long Context Memory Benchmark}.\\
2024.

\bibitem{wu2025longmem}
Z. Wu, et al.\\
\textit{LongMemEval: Evaluating Long-Term Memory in LLM Agents}.\\
2025.

\bibitem{hbit}
H-Bit Protocol.\\
\textit{Universal Cryptographic Steganography System}.\\
2026.

\end{thebibliography}

\end{document}
